\begin{abstract}
We solve the Kervaire invariant problem at the prime three for the standard Kervaire classes in the mod-three Adams spectral sequence. These classes survive precisely at indices $0$, $2$, and $3$, and each surviving detection coset contains an element of additive order three. In particular, there is an order-three Kervaire element in stem $322$.
\end{abstract}

\section{Introduction}\label{sec:introduction}

The odd-primary Kervaire invariant problem asks which distinguished classes on the Adams two-line represent stable homotopy classes of spheres. Its strong form asks for representatives killed by the prime. These are different questions: permanence determines an associated-graded class, while the order of a representative also depends on an extension in actual homotopy. We solve both questions for the standard family at the prime three.

Write
\[
 E_2^{s,t}=\Ext_{\calA_3}^{s,t}(\F_3,\F_3)
 \Longrightarrow (\pi^S_{t-s})^{\wedge}_3,
 \qquad
 b_j\in E_2^{2,\,4\cdot3^{j+1}}\quad(j\geq0),
\]
where $\calA_3$ is the mod-three Steenrod algebra and $b_j$ is the standard Kervaire class. The spectral-sequence and family conventions are those of \cite[Theorems~2.1.1 and~3.4.4]{Ravenel}. Define
\begin{align*}
 K&=\{j\geq0:b_j\text{ survives nontrivially to }E_\infty\},\\
 K_3&=\{j\in K:\text{some class detected by }b_j
                      \text{ has additive order }3\}.
\end{align*}
The second definition concerns existence within a detection coset. It does not require every representative of that coset to have the same additive order.

\begin{theorem}\label{thm:classification}
At the prime three,
\[
                       K=K_3=\{0,2,3\}.
\]
Thus the surviving classes occur in stems $10$, $106$, and $322$, and each admits an additive-order-three representative. The class $b_1$ and every $b_j$ with $j\geq4$ fail to survive.
\end{theorem}

Positive stable stems are finite \cite[Theorem~1.1.8]{Ravenel}. Consequently, the assertion is unchanged if the sphere is taken to be $3$-local instead of $3$-complete. We use the local sphere for the Moore-spectrum construction and complete coefficient rings for Morava $E$-theory. The problem here is the standard odd-primary Adams-family problem; the classical prime-two Kervaire invariant one problem has a different family and dimension convention.

\subsection{Unstable consequences}

Throughout this subsection all spaces, maps, fibres, and homotopy
equivalences are localized at three. For \(n\geq1\), write
\(S^{2n+1}\{3\}\) for the homotopy fibre of the degree-three self-map of
\(S^{2n+1}\), and write
\(E^2:S^{2n-1}\to\Omega^2S^{2n+1}\) for the double suspension.
In Amelotte's notation, the Anick H-space \(T^{2n+1}(3)\) fits into a
fibration sequence
\[
 \Omega^2S^{2n+1}\xrightarrow{\varphi}S^{2n-1}
 \longrightarrow T^{2n+1}(3)\longrightarrow\Omega S^{2n+1},
\]
where \(E^2\circ\varphi\) is homotopic to the third-power map on
\(\Omega^2S^{2n+1}\) \cite[Section~1]{Amelotte}.
Let \(W_n\) be the homotopy fibre of \(E^2\), and let \(BW_n\) be
Gray's classifying space, so that \(\Omega BW_n\simeq W_n\).

\begin{corollary}\label{intro:unstable-consequences}
\begin{enumerate}[label=\textup{(\roman*)}]
\item There are homotopy equivalences of H-spaces
\[
 \Omega S^{163}\{3\}\simeq T^{163}(3)\times\Omega T^{487}(3),
 \qquad BW_{81}\simeq\Omega T^{487}(3).
\]
Consequently, $W_{81}\simeq\Omega^2T^{487}(3)$.
\item For every integer $n>1$, the underlying space $\Omega S^{2n+1}\{3\}$ has a nontrivial product decomposition if and only if $n\in\{3,27,81\}$. For every other $n>1$ it is mod-three atomic, hence indecomposable. Here a product is nontrivial when both factors are noncontractible.
\item For every integer $n>1$, the homotopy type $T^{2n+1}(3)$ admits a homotopy-associative H-space multiplication if and only if $n\in\{3,27,81\}$. The multiplication on $T^{163}(3)$ supplied by the displayed product is homotopy associative and homotopy commutative. For this multiplication, the third-power map on $T^{163}(3)$ is null-homotopic, as is the third-power map on the loop H-space $\Omega T^{487}(3)$.
\end{enumerate}
\end{corollary}

\begin{proof}
For $n=3^r$ with $r\geq1$, the strong Kervaire condition in stem $4n-2$ asks for an order-three representative detected by $b_{r-1}$; such a representative is Amelotte's $\theta_r$. This is the index change in Lemma~\ref{cube:strong-criterion}. Theorem~\ref{thm:classification} therefore gives this condition exactly for $r=1,3,4$, or $n=3,27,81$. Amelotte's Theorem~1.1 \cite{Amelotte} supplies a nontrivial H-space product at each of these indices, with Anick dimensions $2\cdot3^r+1$ and $2\cdot3^{r+1}+1$. At $r=4$ these dimensions are $163$ and $487$, and the same theorem gives the displayed $BW_{81}$ equivalence. Looping that equivalence gives the assertion about $W_{81}$.

In the $n>1$ setting of Selick's Theorem~3.3, his Corollary~3.4 \cite{Selick} makes $\Omega S^{2n+1}\{3\}$ mod-three atomic when the strong Kervaire condition in stem $4n-2$ fails; this includes every $n$ that is not a power of three. For the remaining powers, Theorem~\ref{thm:classification} excludes $r=2$ and all $r\geq5$. Atomicity excludes any product with two noncontractible factors, not only a product of the displayed Anick form; see \cite[Definition~3.1 and its discussion]{Selick}. This proves \textup{(ii)}. Amelotte's Theorem~1.4 and its proof \cite{Amelotte} identify, for every $n>1$, the existence of a homotopy-associative multiplication with the same order-three strong Kervaire condition in stem $4n-2$; its necessity also forces $n=3^r$. This proves the classification in \textup{(iii)}.

Finally, the proof of Amelotte's Corollary~3.3 \cite{Amelotte} uses that the loop H-space $\Omega S^{2n+1}\{3\}$ is homotopy associative and homotopy commutative and has null third-power map. These properties pass to H-space retracts. Repeating that argument for the $r=4$ product proves the stated commutativity and power-map assertions. The multiplication assertion is existential, and no conclusion about power maps or loop-space structures on other $BW_n$ is asserted at the remaining indices.
\end{proof}

\subsection{Prior work and the prime-three obstruction}

The name reflects an analogy on the Adams two-line. Browder identified the classical framed Kervaire invariant one problem with the permanence of the squares $h_j^2$ \cite[Theorem~7.1]{Browder}; at an odd prime, the distinguished classes $b_j$ are associated with Adem relations among reduced powers \cite[Section~1]{Ravenel1978}. Toda proved the first exclusion in the odd-primary family: $b_1$ supports a nonzero differential at every odd prime \cite{TodaRelation}.

The Adams--Novikov spectral sequence brings the formal-group structure of complex cobordism into the calculation of stable stems. Quillen identified its universal formal group and the $p$-typical Brown--Peterson summand \cite{Quillen}; the one- and two-line calculations of Miller, Ravenel, and Wilson then described the alpha and beta families and their images under Thom reduction \cite{MRW}. Ravenel's proof of odd-primary nonexistence uses Morava's formal-group viewpoint and restriction to a cyclic subgroup of a stabilizer group. For every odd prime $p\geq5$, it excludes all standard Kervaire classes $b_j$ with $j\geq1$ \cite[Theorem~17]{Ravenel1978}.

At three, the same conclusion fails. Toda's classical fifth differential kills $\beta_{3/3}$, but the next candidate encounters a different phenomenon: the fate of one Adams--Novikov beta representative need not determine the fate of its mod-three Adams image. The classical survivor in stem $106$ is detected by a corrected combination of $\beta_{9/9}$ and $\beta_7$, although $\beta_{9/9}$ itself supports a differential. This distinction is already present in Ravenel's discussion of the prime-three case \cite{Ravenel1978}; see also \cite[Theorem~3.4.4(c) and the discussion after Proposition~6.4.11]{Ravenel}. We use the beta-family and Thom-reduction calculations of Miller--Ravenel--Wilson \cite[Theorems~2.2, 2.6, and~9.4]{MRW}, together with the classical fifth and ninth differentials in the form recorded by Belmont--Shimomura \cite[Lemma~2.5]{BelmontShimomura}. Our detection argument retains every two-line preimage of the Adams class, including these corrections.

The order refinement has its own geometric meaning. Selick's strong odd-primary criterion identifies an order-$p$ Kervaire representative with a stable Moore-to-sphere map whose cone has a specified nonzero primary reduced power \cite{Selick}. Amelotte relates this criterion to product decompositions of looped degree-$p$ sphere fibres and records the order-three representative in stem $106$ \cite[Theorem~1.1 and the proof of Corollary~1.2]{Amelotte}. His class $\theta_r$ is detected by $b_{r-1}$; thus $b_2$ here corresponds to $\theta_3$, and the new $b_3$ class corresponds to $\theta_4$. In the positive branch of our proof, the Moore cofibre sequence is retained through the final correction, so it supplies the additive order in actual sphere homotopy.

An antecedent for the subgroup calculation appears in Hill's report on joint work with Hopkins and Ravenel \cite[Conjecture~1 and Propositions~2--3]{Hill2007}. Assuming Hopkins's coefficient-action conjecture, that report sketches the $C_p$ homotopy fixed point calculation at heights $(p-1)f$ for $f\geq1$, using formal-module deformations and norm classes. Section~\ref{sec:subgroup} establishes the finite-page statements needed here from the completed coefficient model, fixed geometry, and Euler bound proved in the preceding sections.

Hill, Hopkins, and Ravenel proposed using the action of $C_9$ on height-six Morava $E$-theory at three \cite{HHR2011}. Their published norm and geometric-fixed-point machinery \cite{HHR2016} supplies tools for this approach, separately from its prime-two application. Subsequent odd-primary developments include the $\mu_p$ orientations at height $p-1$ of Hahn--Senger--Wilson \cite{HahnSengerWilson}, and the combination of skeletal Borel and Adams--Novikov filtrations in Angelini-Knoll--Behrens--Belmont--Kong \cite{AKBBK}. These provide context for the equivariant and filtered methods; neither gives a height-six prime-three classification.

For the height-six action itself, Belmont--Ray prove detection \cite[Proposition~0.2]{BR} and, assuming an explicit coefficient-action model, calculate the relevant positive-degree group cohomology \cite[Theorem~0.6]{BR}. Their detection follows the cyclotomic formal-module evaluation in the proposed Hill--Hopkins--Ravenel program. We use that mechanism with an integral target and a finite-layer argument for the image of an actual sphere class. Section~\ref{sec:coefficients} proves the specific completed coefficient model proposed in \cite[Conjecture~0.3]{BR}, including the nonreduced fixed geometry needed for the differential constraints.

The proposed route also included $972$-periodicity and a literal vanishing assertion in degree $-2$ \cite[Conjecture~0.1]{BR}. The latter assertion is false. Proposition~\ref{det:full-gap} proves that an additive semilinear trace lifts through the actual homotopy transfer to an infinite-order element of $\pi_{2n}(E_6^{hC_9})$ for every integer $n$. This refutes the full-gap clause of that conjecture. It leaves the separate periodicity and torsion questions untouched. Theorem~\ref{thm:classification} instead uses finite-page outgoing differentials and requires neither of those assertions.

\subsection{Ideas and proof structure}

The proof has two independent branches beyond the classical endpoints. One excludes all indices at least four by examining their images in the height-six detector. The other constructs the stem-$322$ class, starting from the known order-three class in stem $106$.

For the first branch, let $E=E_6$ be height-six Morava $E$-theory, let $G=C_9$ be its chosen stabilizer subgroup, and let $H=C_3\subset G$. Section~\ref{sec:detector} first proves that an actual sphere class detected by $b_j$, $j\geq2$, has a nonzero finite-order image in $E^{hG}$. It also uses a coefficient of a strict coordinate change, a genuine equivariant norm, and geometric fixed points to prove finite Euler-class bounds for both groups. These give deadlines by which the unit must be hit in the Euler-localized spectral sequences. No infinite Tate convergence assertion is needed.

The coefficient geometry makes those deadlines effective. Section~\ref{sec:coefficients} constructs an actual trace-zero period whose $G$-orbit presents the complete coefficient ring. The fixed deformation problem is a strict formal-module problem. Its Hasse sections, centralizer action, and determinant character constrain equivariant maps between the scalar and exterior parts of coefficient cohomology. These constraints exclude maps that could otherwise occur as differentials. Section~\ref{sec:subgroup} combines those restrictions with classical differentials and the finite Euler bound to force the needed $H$-differentials. Section~\ref{sec:ascent} then retains the integral restriction kernels while passing to $G$.

The final exclusion depends on the direction of a differential. An actual sphere class can acquire an incoming differential after passage from ordinary homotopy fixed points to Tate localization. Vanishing of the localized pages alone would therefore not exclude it. We show instead that the detector for every $j\geq4$ must support a nonzero outgoing differential. A finite array of one-dimensional groups records all remaining degree- and character-compatible differentials. Its translation units are proved to survive before they are used. The known $b_2$ survivor fixes an incoming pairing in this array; that pairing forces an outgoing differential on the entire lattice containing the higher detectors.

The positive branch adapts Toda's cone construction for extended powers \cite[Section~1, Theorem~2]{TodaExtended}, within the broader framework of primary operations developed by May \cite[Sections~2--3]{MayOperations}, and uses Selick's primary-operation criterion. Section~\ref{sec:cube} forms a truncated extended cube of a Moore-spectrum map representing the order-three stem-$106$ class. The cube maps into a four-cell spectrum $T$ containing the Moore spectrum $M=S/3$, with quotient $\Sigma^4M$, and has a nonzero $\mathcal P^{81}$ in its cone. Vanishing of two obstruction components produces a map of Adams--Novikov filtration at least two with the same image in that quotient. Subtracting it gives a map that lifts to $M$ and preserves the primary operation. The remaining actual Moore boundary is the required sphere element killed by three.

Section~\ref{sec:finite} supplies precisely the finite Adams--Novikov information needed for those corrections. It proves a cubic boundary, a zero multiplication map on the sixth page, and vanishing on the tenth page in the higher obstruction filtrations. The fifth differential is constrained by affine linear equations, then by linear consequences of its square-zero identity. Universal kernel and image bounds apply to every differential satisfying those necessary conditions; the argument does not select an unproved differential from that family. A weight estimate covers every remaining higher filtration. Appendix~\ref{app:recurrence} gives the integral structure formulas, terminating resolution, cobar comparison, signed products, and elimination procedures defining all finite tables.

\subsection{Reduction to the main arguments}

We record the classical endpoints, including the order convention, before organizing the remaining proof.

\begin{lemma}\label{intro:classical}
One has $0,2\in K_3$ and $1\notin K$.
\end{lemma}

\begin{proof}
The class $\beta_1\in\pi^S_{10}$ has order three and is detected by $b_0$. The classical stem-$106$ Kervaire element has an order-three representative detected by $b_2$; see \cite[Theorem~3.4.4(c), discussion after Proposition~6.4.11, and Table~A3.4]{Ravenel} and \cite[Theorem~1.1 and Corollary~1.2]{Amelotte}. This is the surviving corrected beta combination, not an assertion of survival for $\beta_{9/9}$ alone.

If $b_1$ survived, Lemma~\ref{det:filtration-comparison} would give a permanent Adams--Novikov two-line class in internal degree $36$ with that Thom reduction. The Miller--Ravenel--Wilson two-line calculation~\cite[Theorem~2.6]{MRW} identifies this group with the order-three group generated by $\beta_{3/3}$; see also \cite[Lemma~6.4.12(i)]{Ravenel}. The classical Toda differential
\[
 d_5(\beta_{3/3})=\pm\alpha_1\beta_1^3\neq0
\]
excludes such a permanent class; see \cite[Lemma~2.5]{BelmontShimomura}.
\end{proof}

\begin{proof}[Proof of Theorem~\ref{thm:classification}]
Lemma~\ref{intro:classical} settles the indices $0$, $1$, and $2$. Theorem~\ref{g:tail-exclusion}, proved in Section~\ref{sec:ascent}, excludes every index $j\geq4$. Theorem~\ref{cube:order-three}, proved in Section~\ref{sec:cube} using the finite algebra of Section~\ref{sec:finite}, constructs an exact-order-three representative detected by $b_3$. These assertions exhaust all nonnegative indices. Since $K_3\subseteq K$ by definition, both sets are as stated.
\end{proof}

\subsection{Conventions and organization}

Spectral-sequence degrees are written $(s,t)$, with differentials of degree $(r,r-1)$, unless a rescaled weight is explicitly introduced. All homotopy-theoretic constructions are at the prime three. The word ``Tate'' refers to Euler localization; it does not mean rationalization. We distinguish the algebraic spectral sequences used to calculate group cohomology from the topological spectral sequences whose differentials constrain homotopy.

Sections~\ref{sec:detector}--\ref{sec:ascent} use $G=C_9$ and its subgroup $H=C_3$. Section~\ref{sec:cube} begins the independent stable Moore-spectrum argument. Section~\ref{sec:finite} resets the algebraic notation and uses $(f,d)=(s,t/4)$. Every invocation of a Laurent generator as a translation is restricted to pages on which it and its inverse are cycles.
